\documentclass[11pt]{article}

    \usepackage[breakable]{tcolorbox}
    \usepackage{parskip} % Stop auto-indenting (to mimic markdown behaviour)
    

    % Basic figure setup, for now with no caption control since it's done
    % automatically by Pandoc (which extracts ![](path) syntax from Markdown).
    \usepackage{graphicx}
    % Maintain compatibility with old templates. Remove in nbconvert 6.0
    \let\Oldincludegraphics\includegraphics
    % Ensure that by default, figures have no caption (until we provide a
    % proper Figure object with a Caption API and a way to capture that
    % in the conversion process - todo).
    \usepackage{caption}
    \DeclareCaptionFormat{nocaption}{}
    \captionsetup{format=nocaption,aboveskip=0pt,belowskip=0pt}

    \usepackage{float}
    \floatplacement{figure}{H} % forces figures to be placed at the correct location
    \usepackage{xcolor} % Allow colors to be defined
    \usepackage{enumerate} % Needed for markdown enumerations to work
    \usepackage{geometry} % Used to adjust the document margins
    \usepackage{amsmath} % Equations
    \usepackage{amssymb} % Equations
    \usepackage{textcomp} % defines textquotesingle
    % Hack from http://tex.stackexchange.com/a/47451/13684:
    \AtBeginDocument{%
        \def\PYZsq{\textquotesingle}% Upright quotes in Pygmentized code
    }
    \usepackage{upquote} % Upright quotes for verbatim code
    \usepackage{eurosym} % defines \euro

    \usepackage{iftex}
    \ifPDFTeX
        \usepackage[T1]{fontenc}
        \IfFileExists{alphabeta.sty}{
              \usepackage{alphabeta}
          }{
              \usepackage[mathletters]{ucs}
              \usepackage[utf8x]{inputenc}
          }
    \else
        \usepackage{fontspec}
        \usepackage{unicode-math}
    \fi

    \usepackage{fancyvrb} % verbatim replacement that allows latex
    \usepackage{grffile} % extends the file name processing of package graphics
                         % to support a larger range
    \makeatletter % fix for old versions of grffile with XeLaTeX
    \@ifpackagelater{grffile}{2019/11/01}
    {
      % Do nothing on new versions
    }
    {
      \def\Gread@@xetex#1{%
        \IfFileExists{"\Gin@base".bb}%
        {\Gread@eps{\Gin@base.bb}}%
        {\Gread@@xetex@aux#1}%
      }
    }
    \makeatother
    \usepackage[Export]{adjustbox} % Used to constrain images to a maximum size
    \adjustboxset{max size={0.9\linewidth}{0.9\paperheight}}

    % The hyperref package gives us a pdf with properly built
    % internal navigation ('pdf bookmarks' for the table of contents,
    % internal cross-reference links, web links for URLs, etc.)
    \usepackage{hyperref}
    % The default LaTeX title has an obnoxious amount of whitespace. By default,
    % titling removes some of it. It also provides customization options.
    \usepackage{titling}
    \usepackage{longtable} % longtable support required by pandoc >1.10
    \usepackage{booktabs}  % table support for pandoc > 1.12.2
    \usepackage{array}     % table support for pandoc >= 2.11.3
    \usepackage{calc}      % table minipage width calculation for pandoc >= 2.11.1
    \usepackage[inline]{enumitem} % IRkernel/repr support (it uses the enumerate* environment)
    \usepackage[normalem]{ulem} % ulem is needed to support strikethroughs (\sout)
                                % normalem makes italics be italics, not underlines
    \usepackage{soul}      % strikethrough (\st) support for pandoc >= 3.0.0
    \usepackage{mathrsfs}
    

    
    % Colors for the hyperref package
    \definecolor{urlcolor}{rgb}{0,.145,.698}
    \definecolor{linkcolor}{rgb}{.71,0.21,0.01}
    \definecolor{citecolor}{rgb}{.12,.54,.11}

    % ANSI colors
    \definecolor{ansi-black}{HTML}{3E424D}
    \definecolor{ansi-black-intense}{HTML}{282C36}
    \definecolor{ansi-red}{HTML}{E75C58}
    \definecolor{ansi-red-intense}{HTML}{B22B31}
    \definecolor{ansi-green}{HTML}{00A250}
    \definecolor{ansi-green-intense}{HTML}{007427}
    \definecolor{ansi-yellow}{HTML}{DDB62B}
    \definecolor{ansi-yellow-intense}{HTML}{B27D12}
    \definecolor{ansi-blue}{HTML}{208FFB}
    \definecolor{ansi-blue-intense}{HTML}{0065CA}
    \definecolor{ansi-magenta}{HTML}{D160C4}
    \definecolor{ansi-magenta-intense}{HTML}{A03196}
    \definecolor{ansi-cyan}{HTML}{60C6C8}
    \definecolor{ansi-cyan-intense}{HTML}{258F8F}
    \definecolor{ansi-white}{HTML}{C5C1B4}
    \definecolor{ansi-white-intense}{HTML}{A1A6B2}
    \definecolor{ansi-default-inverse-fg}{HTML}{FFFFFF}
    \definecolor{ansi-default-inverse-bg}{HTML}{000000}

    % common color for the border for error outputs.
    \definecolor{outerrorbackground}{HTML}{FFDFDF}

    % commands and environments needed by pandoc snippets
    % extracted from the output of `pandoc -s`
    \providecommand{\tightlist}{%
      \setlength{\itemsep}{0pt}\setlength{\parskip}{0pt}}
    \DefineVerbatimEnvironment{Highlighting}{Verbatim}{commandchars=\\\{\}}
    % Add ',fontsize=\small' for more characters per line
    \newenvironment{Shaded}{}{}
    \newcommand{\KeywordTok}[1]{\textcolor[rgb]{0.00,0.44,0.13}{\textbf{{#1}}}}
    \newcommand{\DataTypeTok}[1]{\textcolor[rgb]{0.56,0.13,0.00}{{#1}}}
    \newcommand{\DecValTok}[1]{\textcolor[rgb]{0.25,0.63,0.44}{{#1}}}
    \newcommand{\BaseNTok}[1]{\textcolor[rgb]{0.25,0.63,0.44}{{#1}}}
    \newcommand{\FloatTok}[1]{\textcolor[rgb]{0.25,0.63,0.44}{{#1}}}
    \newcommand{\CharTok}[1]{\textcolor[rgb]{0.25,0.44,0.63}{{#1}}}
    \newcommand{\StringTok}[1]{\textcolor[rgb]{0.25,0.44,0.63}{{#1}}}
    \newcommand{\CommentTok}[1]{\textcolor[rgb]{0.38,0.63,0.69}{\textit{{#1}}}}
    \newcommand{\OtherTok}[1]{\textcolor[rgb]{0.00,0.44,0.13}{{#1}}}
    \newcommand{\AlertTok}[1]{\textcolor[rgb]{1.00,0.00,0.00}{\textbf{{#1}}}}
    \newcommand{\FunctionTok}[1]{\textcolor[rgb]{0.02,0.16,0.49}{{#1}}}
    \newcommand{\RegionMarkerTok}[1]{{#1}}
    \newcommand{\ErrorTok}[1]{\textcolor[rgb]{1.00,0.00,0.00}{\textbf{{#1}}}}
    \newcommand{\NormalTok}[1]{{#1}}

    % Additional commands for more recent versions of Pandoc
    \newcommand{\ConstantTok}[1]{\textcolor[rgb]{0.53,0.00,0.00}{{#1}}}
    \newcommand{\SpecialCharTok}[1]{\textcolor[rgb]{0.25,0.44,0.63}{{#1}}}
    \newcommand{\VerbatimStringTok}[1]{\textcolor[rgb]{0.25,0.44,0.63}{{#1}}}
    \newcommand{\SpecialStringTok}[1]{\textcolor[rgb]{0.73,0.40,0.53}{{#1}}}
    \newcommand{\ImportTok}[1]{{#1}}
    \newcommand{\DocumentationTok}[1]{\textcolor[rgb]{0.73,0.13,0.13}{\textit{{#1}}}}
    \newcommand{\AnnotationTok}[1]{\textcolor[rgb]{0.38,0.63,0.69}{\textbf{\textit{{#1}}}}}
    \newcommand{\CommentVarTok}[1]{\textcolor[rgb]{0.38,0.63,0.69}{\textbf{\textit{{#1}}}}}
    \newcommand{\VariableTok}[1]{\textcolor[rgb]{0.10,0.09,0.49}{{#1}}}
    \newcommand{\ControlFlowTok}[1]{\textcolor[rgb]{0.00,0.44,0.13}{\textbf{{#1}}}}
    \newcommand{\OperatorTok}[1]{\textcolor[rgb]{0.40,0.40,0.40}{{#1}}}
    \newcommand{\BuiltInTok}[1]{{#1}}
    \newcommand{\ExtensionTok}[1]{{#1}}
    \newcommand{\PreprocessorTok}[1]{\textcolor[rgb]{0.74,0.48,0.00}{{#1}}}
    \newcommand{\AttributeTok}[1]{\textcolor[rgb]{0.49,0.56,0.16}{{#1}}}
    \newcommand{\InformationTok}[1]{\textcolor[rgb]{0.38,0.63,0.69}{\textbf{\textit{{#1}}}}}
    \newcommand{\WarningTok}[1]{\textcolor[rgb]{0.38,0.63,0.69}{\textbf{\textit{{#1}}}}}


    % Define a nice break command that doesn't care if a line doesn't already
    % exist.
    \def\br{\hspace*{\fill} \\* }
    % Math Jax compatibility definitions
    \def\gt{>}
    \def\lt{<}
    \let\Oldtex\TeX
    \let\Oldlatex\LaTeX
    \renewcommand{\TeX}{\textrm{\Oldtex}}
    \renewcommand{\LaTeX}{\textrm{\Oldlatex}}
    % Document parameters
    % Document title
    \title{Project2\_Launch\_JansenRinck}
    
    
    
    
    
    
    
% Pygments definitions
\makeatletter
\def\PY@reset{\let\PY@it=\relax \let\PY@bf=\relax%
    \let\PY@ul=\relax \let\PY@tc=\relax%
    \let\PY@bc=\relax \let\PY@ff=\relax}
\def\PY@tok#1{\csname PY@tok@#1\endcsname}
\def\PY@toks#1+{\ifx\relax#1\empty\else%
    \PY@tok{#1}\expandafter\PY@toks\fi}
\def\PY@do#1{\PY@bc{\PY@tc{\PY@ul{%
    \PY@it{\PY@bf{\PY@ff{#1}}}}}}}
\def\PY#1#2{\PY@reset\PY@toks#1+\relax+\PY@do{#2}}

\@namedef{PY@tok@w}{\def\PY@tc##1{\textcolor[rgb]{0.73,0.73,0.73}{##1}}}
\@namedef{PY@tok@c}{\let\PY@it=\textit\def\PY@tc##1{\textcolor[rgb]{0.24,0.48,0.48}{##1}}}
\@namedef{PY@tok@cp}{\def\PY@tc##1{\textcolor[rgb]{0.61,0.40,0.00}{##1}}}
\@namedef{PY@tok@k}{\let\PY@bf=\textbf\def\PY@tc##1{\textcolor[rgb]{0.00,0.50,0.00}{##1}}}
\@namedef{PY@tok@kp}{\def\PY@tc##1{\textcolor[rgb]{0.00,0.50,0.00}{##1}}}
\@namedef{PY@tok@kt}{\def\PY@tc##1{\textcolor[rgb]{0.69,0.00,0.25}{##1}}}
\@namedef{PY@tok@o}{\def\PY@tc##1{\textcolor[rgb]{0.40,0.40,0.40}{##1}}}
\@namedef{PY@tok@ow}{\let\PY@bf=\textbf\def\PY@tc##1{\textcolor[rgb]{0.67,0.13,1.00}{##1}}}
\@namedef{PY@tok@nb}{\def\PY@tc##1{\textcolor[rgb]{0.00,0.50,0.00}{##1}}}
\@namedef{PY@tok@nf}{\def\PY@tc##1{\textcolor[rgb]{0.00,0.00,1.00}{##1}}}
\@namedef{PY@tok@nc}{\let\PY@bf=\textbf\def\PY@tc##1{\textcolor[rgb]{0.00,0.00,1.00}{##1}}}
\@namedef{PY@tok@nn}{\let\PY@bf=\textbf\def\PY@tc##1{\textcolor[rgb]{0.00,0.00,1.00}{##1}}}
\@namedef{PY@tok@ne}{\let\PY@bf=\textbf\def\PY@tc##1{\textcolor[rgb]{0.80,0.25,0.22}{##1}}}
\@namedef{PY@tok@nv}{\def\PY@tc##1{\textcolor[rgb]{0.10,0.09,0.49}{##1}}}
\@namedef{PY@tok@no}{\def\PY@tc##1{\textcolor[rgb]{0.53,0.00,0.00}{##1}}}
\@namedef{PY@tok@nl}{\def\PY@tc##1{\textcolor[rgb]{0.46,0.46,0.00}{##1}}}
\@namedef{PY@tok@ni}{\let\PY@bf=\textbf\def\PY@tc##1{\textcolor[rgb]{0.44,0.44,0.44}{##1}}}
\@namedef{PY@tok@na}{\def\PY@tc##1{\textcolor[rgb]{0.41,0.47,0.13}{##1}}}
\@namedef{PY@tok@nt}{\let\PY@bf=\textbf\def\PY@tc##1{\textcolor[rgb]{0.00,0.50,0.00}{##1}}}
\@namedef{PY@tok@nd}{\def\PY@tc##1{\textcolor[rgb]{0.67,0.13,1.00}{##1}}}
\@namedef{PY@tok@s}{\def\PY@tc##1{\textcolor[rgb]{0.73,0.13,0.13}{##1}}}
\@namedef{PY@tok@sd}{\let\PY@it=\textit\def\PY@tc##1{\textcolor[rgb]{0.73,0.13,0.13}{##1}}}
\@namedef{PY@tok@si}{\let\PY@bf=\textbf\def\PY@tc##1{\textcolor[rgb]{0.64,0.35,0.47}{##1}}}
\@namedef{PY@tok@se}{\let\PY@bf=\textbf\def\PY@tc##1{\textcolor[rgb]{0.67,0.36,0.12}{##1}}}
\@namedef{PY@tok@sr}{\def\PY@tc##1{\textcolor[rgb]{0.64,0.35,0.47}{##1}}}
\@namedef{PY@tok@ss}{\def\PY@tc##1{\textcolor[rgb]{0.10,0.09,0.49}{##1}}}
\@namedef{PY@tok@sx}{\def\PY@tc##1{\textcolor[rgb]{0.00,0.50,0.00}{##1}}}
\@namedef{PY@tok@m}{\def\PY@tc##1{\textcolor[rgb]{0.40,0.40,0.40}{##1}}}
\@namedef{PY@tok@gh}{\let\PY@bf=\textbf\def\PY@tc##1{\textcolor[rgb]{0.00,0.00,0.50}{##1}}}
\@namedef{PY@tok@gu}{\let\PY@bf=\textbf\def\PY@tc##1{\textcolor[rgb]{0.50,0.00,0.50}{##1}}}
\@namedef{PY@tok@gd}{\def\PY@tc##1{\textcolor[rgb]{0.63,0.00,0.00}{##1}}}
\@namedef{PY@tok@gi}{\def\PY@tc##1{\textcolor[rgb]{0.00,0.52,0.00}{##1}}}
\@namedef{PY@tok@gr}{\def\PY@tc##1{\textcolor[rgb]{0.89,0.00,0.00}{##1}}}
\@namedef{PY@tok@ge}{\let\PY@it=\textit}
\@namedef{PY@tok@gs}{\let\PY@bf=\textbf}
\@namedef{PY@tok@ges}{\let\PY@bf=\textbf\let\PY@it=\textit}
\@namedef{PY@tok@gp}{\let\PY@bf=\textbf\def\PY@tc##1{\textcolor[rgb]{0.00,0.00,0.50}{##1}}}
\@namedef{PY@tok@go}{\def\PY@tc##1{\textcolor[rgb]{0.44,0.44,0.44}{##1}}}
\@namedef{PY@tok@gt}{\def\PY@tc##1{\textcolor[rgb]{0.00,0.27,0.87}{##1}}}
\@namedef{PY@tok@err}{\def\PY@bc##1{{\setlength{\fboxsep}{\string -\fboxrule}\fcolorbox[rgb]{1.00,0.00,0.00}{1,1,1}{\strut ##1}}}}
\@namedef{PY@tok@kc}{\let\PY@bf=\textbf\def\PY@tc##1{\textcolor[rgb]{0.00,0.50,0.00}{##1}}}
\@namedef{PY@tok@kd}{\let\PY@bf=\textbf\def\PY@tc##1{\textcolor[rgb]{0.00,0.50,0.00}{##1}}}
\@namedef{PY@tok@kn}{\let\PY@bf=\textbf\def\PY@tc##1{\textcolor[rgb]{0.00,0.50,0.00}{##1}}}
\@namedef{PY@tok@kr}{\let\PY@bf=\textbf\def\PY@tc##1{\textcolor[rgb]{0.00,0.50,0.00}{##1}}}
\@namedef{PY@tok@bp}{\def\PY@tc##1{\textcolor[rgb]{0.00,0.50,0.00}{##1}}}
\@namedef{PY@tok@fm}{\def\PY@tc##1{\textcolor[rgb]{0.00,0.00,1.00}{##1}}}
\@namedef{PY@tok@vc}{\def\PY@tc##1{\textcolor[rgb]{0.10,0.09,0.49}{##1}}}
\@namedef{PY@tok@vg}{\def\PY@tc##1{\textcolor[rgb]{0.10,0.09,0.49}{##1}}}
\@namedef{PY@tok@vi}{\def\PY@tc##1{\textcolor[rgb]{0.10,0.09,0.49}{##1}}}
\@namedef{PY@tok@vm}{\def\PY@tc##1{\textcolor[rgb]{0.10,0.09,0.49}{##1}}}
\@namedef{PY@tok@sa}{\def\PY@tc##1{\textcolor[rgb]{0.73,0.13,0.13}{##1}}}
\@namedef{PY@tok@sb}{\def\PY@tc##1{\textcolor[rgb]{0.73,0.13,0.13}{##1}}}
\@namedef{PY@tok@sc}{\def\PY@tc##1{\textcolor[rgb]{0.73,0.13,0.13}{##1}}}
\@namedef{PY@tok@dl}{\def\PY@tc##1{\textcolor[rgb]{0.73,0.13,0.13}{##1}}}
\@namedef{PY@tok@s2}{\def\PY@tc##1{\textcolor[rgb]{0.73,0.13,0.13}{##1}}}
\@namedef{PY@tok@sh}{\def\PY@tc##1{\textcolor[rgb]{0.73,0.13,0.13}{##1}}}
\@namedef{PY@tok@s1}{\def\PY@tc##1{\textcolor[rgb]{0.73,0.13,0.13}{##1}}}
\@namedef{PY@tok@mb}{\def\PY@tc##1{\textcolor[rgb]{0.40,0.40,0.40}{##1}}}
\@namedef{PY@tok@mf}{\def\PY@tc##1{\textcolor[rgb]{0.40,0.40,0.40}{##1}}}
\@namedef{PY@tok@mh}{\def\PY@tc##1{\textcolor[rgb]{0.40,0.40,0.40}{##1}}}
\@namedef{PY@tok@mi}{\def\PY@tc##1{\textcolor[rgb]{0.40,0.40,0.40}{##1}}}
\@namedef{PY@tok@il}{\def\PY@tc##1{\textcolor[rgb]{0.40,0.40,0.40}{##1}}}
\@namedef{PY@tok@mo}{\def\PY@tc##1{\textcolor[rgb]{0.40,0.40,0.40}{##1}}}
\@namedef{PY@tok@ch}{\let\PY@it=\textit\def\PY@tc##1{\textcolor[rgb]{0.24,0.48,0.48}{##1}}}
\@namedef{PY@tok@cm}{\let\PY@it=\textit\def\PY@tc##1{\textcolor[rgb]{0.24,0.48,0.48}{##1}}}
\@namedef{PY@tok@cpf}{\let\PY@it=\textit\def\PY@tc##1{\textcolor[rgb]{0.24,0.48,0.48}{##1}}}
\@namedef{PY@tok@c1}{\let\PY@it=\textit\def\PY@tc##1{\textcolor[rgb]{0.24,0.48,0.48}{##1}}}
\@namedef{PY@tok@cs}{\let\PY@it=\textit\def\PY@tc##1{\textcolor[rgb]{0.24,0.48,0.48}{##1}}}

\def\PYZbs{\char`\\}
\def\PYZus{\char`\_}
\def\PYZob{\char`\{}
\def\PYZcb{\char`\}}
\def\PYZca{\char`\^}
\def\PYZam{\char`\&}
\def\PYZlt{\char`\<}
\def\PYZgt{\char`\>}
\def\PYZsh{\char`\#}
\def\PYZpc{\char`\%}
\def\PYZdl{\char`\$}
\def\PYZhy{\char`\-}
\def\PYZsq{\char`\'}
\def\PYZdq{\char`\"}
\def\PYZti{\char`\~}
% for compatibility with earlier versions
\def\PYZat{@}
\def\PYZlb{[}
\def\PYZrb{]}
\makeatother


    % For linebreaks inside Verbatim environment from package fancyvrb.
    \makeatletter
        \newbox\Wrappedcontinuationbox
        \newbox\Wrappedvisiblespacebox
        \newcommand*\Wrappedvisiblespace {\textcolor{red}{\textvisiblespace}}
        \newcommand*\Wrappedcontinuationsymbol {\textcolor{red}{\llap{\tiny$\m@th\hookrightarrow$}}}
        \newcommand*\Wrappedcontinuationindent {3ex }
        \newcommand*\Wrappedafterbreak {\kern\Wrappedcontinuationindent\copy\Wrappedcontinuationbox}
        % Take advantage of the already applied Pygments mark-up to insert
        % potential linebreaks for TeX processing.
        %        {, <, #, %, $, ' and ": go to next line.
        %        _, }, ^, &, >, - and ~: stay at end of broken line.
        % Use of \textquotesingle for straight quote.
        \newcommand*\Wrappedbreaksatspecials {%
            \def\PYGZus{\discretionary{\char`\_}{\Wrappedafterbreak}{\char`\_}}%
            \def\PYGZob{\discretionary{}{\Wrappedafterbreak\char`\{}{\char`\{}}%
            \def\PYGZcb{\discretionary{\char`\}}{\Wrappedafterbreak}{\char`\}}}%
            \def\PYGZca{\discretionary{\char`\^}{\Wrappedafterbreak}{\char`\^}}%
            \def\PYGZam{\discretionary{\char`\&}{\Wrappedafterbreak}{\char`\&}}%
            \def\PYGZlt{\discretionary{}{\Wrappedafterbreak\char`\<}{\char`\<}}%
            \def\PYGZgt{\discretionary{\char`\>}{\Wrappedafterbreak}{\char`\>}}%
            \def\PYGZsh{\discretionary{}{\Wrappedafterbreak\char`\#}{\char`\#}}%
            \def\PYGZpc{\discretionary{}{\Wrappedafterbreak\char`\%}{\char`\%}}%
            \def\PYGZdl{\discretionary{}{\Wrappedafterbreak\char`\$}{\char`\$}}%
            \def\PYGZhy{\discretionary{\char`\-}{\Wrappedafterbreak}{\char`\-}}%
            \def\PYGZsq{\discretionary{}{\Wrappedafterbreak\textquotesingle}{\textquotesingle}}%
            \def\PYGZdq{\discretionary{}{\Wrappedafterbreak\char`\"}{\char`\"}}%
            \def\PYGZti{\discretionary{\char`\~}{\Wrappedafterbreak}{\char`\~}}%
        }
        % Some characters . , ; ? ! / are not pygmentized.
        % This macro makes them "active" and they will insert potential linebreaks
        \newcommand*\Wrappedbreaksatpunct {%
            \lccode`\~`\.\lowercase{\def~}{\discretionary{\hbox{\char`\.}}{\Wrappedafterbreak}{\hbox{\char`\.}}}%
            \lccode`\~`\,\lowercase{\def~}{\discretionary{\hbox{\char`\,}}{\Wrappedafterbreak}{\hbox{\char`\,}}}%
            \lccode`\~`\;\lowercase{\def~}{\discretionary{\hbox{\char`\;}}{\Wrappedafterbreak}{\hbox{\char`\;}}}%
            \lccode`\~`\:\lowercase{\def~}{\discretionary{\hbox{\char`\:}}{\Wrappedafterbreak}{\hbox{\char`\:}}}%
            \lccode`\~`\?\lowercase{\def~}{\discretionary{\hbox{\char`\?}}{\Wrappedafterbreak}{\hbox{\char`\?}}}%
            \lccode`\~`\!\lowercase{\def~}{\discretionary{\hbox{\char`\!}}{\Wrappedafterbreak}{\hbox{\char`\!}}}%
            \lccode`\~`\/\lowercase{\def~}{\discretionary{\hbox{\char`\/}}{\Wrappedafterbreak}{\hbox{\char`\/}}}%
            \catcode`\.\active
            \catcode`\,\active
            \catcode`\;\active
            \catcode`\:\active
            \catcode`\?\active
            \catcode`\!\active
            \catcode`\/\active
            \lccode`\~`\~
        }
    \makeatother

    \let\OriginalVerbatim=\Verbatim
    \makeatletter
    \renewcommand{\Verbatim}[1][1]{%
        %\parskip\z@skip
        \sbox\Wrappedcontinuationbox {\Wrappedcontinuationsymbol}%
        \sbox\Wrappedvisiblespacebox {\FV@SetupFont\Wrappedvisiblespace}%
        \def\FancyVerbFormatLine ##1{\hsize\linewidth
            \vtop{\raggedright\hyphenpenalty\z@\exhyphenpenalty\z@
                \doublehyphendemerits\z@\finalhyphendemerits\z@
                \strut ##1\strut}%
        }%
        % If the linebreak is at a space, the latter will be displayed as visible
        % space at end of first line, and a continuation symbol starts next line.
        % Stretch/shrink are however usually zero for typewriter font.
        \def\FV@Space {%
            \nobreak\hskip\z@ plus\fontdimen3\font minus\fontdimen4\font
            \discretionary{\copy\Wrappedvisiblespacebox}{\Wrappedafterbreak}
            {\kern\fontdimen2\font}%
        }%

        % Allow breaks at special characters using \PYG... macros.
        \Wrappedbreaksatspecials
        % Breaks at punctuation characters . , ; ? ! and / need catcode=\active
        \OriginalVerbatim[#1,codes*=\Wrappedbreaksatpunct]%
    }
    \makeatother

    % Exact colors from NB
    \definecolor{incolor}{HTML}{303F9F}
    \definecolor{outcolor}{HTML}{D84315}
    \definecolor{cellborder}{HTML}{CFCFCF}
    \definecolor{cellbackground}{HTML}{F7F7F7}

    % prompt
    \makeatletter
    \newcommand{\boxspacing}{\kern\kvtcb@left@rule\kern\kvtcb@boxsep}
    \makeatother
    \newcommand{\prompt}[4]{
        {\ttfamily\llap{{\color{#2}[#3]:\hspace{3pt}#4}}\vspace{-\baselineskip}}
    }
    

    
    % Prevent overflowing lines due to hard-to-break entities
    \sloppy
    % Setup hyperref package
    \hypersetup{
      breaklinks=true,  % so long urls are correctly broken across lines
      colorlinks=true,
      urlcolor=urlcolor,
      linkcolor=linkcolor,
      citecolor=citecolor,
      }
    % Slightly bigger margins than the latex defaults
    
    \geometry{verbose,tmargin=0.5in,bmargin=0.5in,lmargin=0.5in,rmargin=0.5in}
    
    

\begin{document}
    
    \maketitle
    
    

    
    \hypertarget{project-2-launch}{%
\section{Project 2: Launch}\label{project-2-launch}}

    \hypertarget{research-question}{%
\subsection{Research Question:}\label{research-question}}

Between 2018 and 2020, does the rate at which lightning strikes convert
into wildfire ignitions vary by fuel type across the ecoregions of the
contiguous United States, independent of raw lightning exposure?

    \hypertarget{data}{%
\section{Data:}\label{data}}

\begin{enumerate}
\def\labelenumi{\arabic{enumi}.}
\tightlist
\item
  \href{https://www.fs.usda.gov/rds/archive/catalog/RDS-2013-0009.6}{Spatial
  wildfire occurrence data for the United States, 1992-2020 (6th
  Edition)(USFS)}

  \begin{itemize}
  \tightlist
  \item
    Federal, State and Local fire data aggreagated
  \item
    Discovery date, final fire size, point location (1 sq mile
    accuracy), wildfire-cause
  \item
    2.3 million geo-referenced wildfire records
  \end{itemize}
\item
  \href{https://search.earthdata.nasa.gov/search/granules/collection-details?p=C3894650367-GHRC_DAAC\&pg\%5B0\%5D\%5Bv\%5D=f\&pg\%5B0\%5D\%5Bgsk\%5D=-start_date\&q=Lightning\%20Mapper\&lat=40.780546668112606\&long=-83.81276234265948\&zoom=3.990854889796647}{Combined
  Geostationary Lightning Mapper (GLM) Hourly Thunder Hour Data Product
  (NASA)}

  \begin{itemize}
  \tightlist
  \item
    2019 - 2025
  \item
    Two sattelites GOES-18 (west) and GOES-19 (east)
  \item
    Thunder hour: an hour with detected lightning at a location - a
    frequency-based storm-activity metric, less biased by a sensor's
    geographic detection-sensitivity variation than raw flash counts
  \item
    Resolution: 5 km
  \end{itemize}
\item
  \href{https://www.sciencebase.gov/catalog/item/6a1602c7b66b013d45990b82}{LANDFIRE's
  40 Scott and Burgan Fire Behavior Fuel Model (FBFM40), CONUS raster
  (USFS \& USGS}

  \begin{itemize}
  \tightlist
  \item
    40 fuel models classifying surface fuel loading by type (grass,
    brush, timber, slash), size class, bed depth, and moisture of
    extinction
  \item
    Fuel assignments aren't static - herbaceous fuels shift between live
    and dead through the growing season, and burned or otherwise
    disturbed areas are updated to reflect estimated recovery through
    2022
  \end{itemize}
\item
  \href{https://www.epa.gov/eco-research/level-iii-and-iv-ecoregions-continental-united-states}{Level
  III Ecoregions of the Conterminous United States (EPA)}

  \begin{itemize}
  \tightlist
  \item
    Ecoregions mark areas of broadly similar ecosystems - a spatial
    framework for managing and comparing environmental resources across
    agencies
  \item
    Delineated from patterns across multiple factors (geology,
    vegetation, climate, soils, land use, wildlife, hydrology), weighted
    differently region to region
  \item
    Four Levels: Level I (15 regions, North America) → Level II (50) →
    Level III (85 in the conterminous US) → Level IV (further
    subdivisions)
  \end{itemize}
\item
  \href{https://in-form-nifc.hub.arcgis.com/pages/in-form-faqs}{Backup
  Wildfire Data: Interagency Fire Occurrence Reporting Modules (InFORM}

  \begin{itemize}
  \tightlist
  \item
    Ongling incorporation of historic data across other agencies is
    still in progress
  \item
    Data still needs to be backfilled
  \end{itemize}
\end{enumerate}

    \hypertarget{data-processing}{%
\subsection{Data Processing:}\label{data-processing}}

\begin{enumerate}
\def\labelenumi{\arabic{enumi}.}
\tightlist
\item
  CRS

  \begin{itemize}
  \tightlist
  \item
    Albers Equal Area Conic

    \begin{itemize}
    \tightlist
    \item
      Preserves area calculation accuracy
    \item
      Minimal distortion for CONUS
    \end{itemize}
  \end{itemize}
\item
  Zonal Statistics at CONUS Scale for LANDFIRE Raster

  \begin{itemize}
  \tightlist
  \item
    Windowed (chunked) reads for the raster data set to avoid memory
    issues in Docker
  \end{itemize}
\item
  Lightning Data Cropping to CONUS
\item
  LANDFIRE Attribute Table

  \begin{itemize}
  \tightlist
  \item
    Translate values to fuel model classes
  \end{itemize}
\item
  Categorize Fuel Types

  \begin{itemize}
  \tightlist
  \item
    Single majority class per eco region?
  \item
    Proportions of fuel class per eco region?
  \item
    Lightning strikes per fuel class
  \end{itemize}
\end{enumerate}

    \hypertarget{analysis-methods}{%
\subsection{Analysis Methods:}\label{analysis-methods}}

\hypertarget{spatial-joins-and-overlays}{%
\subsubsection{Spatial joins and
overlays}\label{spatial-joins-and-overlays}}

\begin{itemize}
\tightlist
\item
  Fire points joined to ecoregion polygons (point-in-polygon)
\item
  Fire points sampled against the fuel raster (tagging each ignition
  with its fuel class)
\item
  Lightning grid cells overlaid on the fuel raster to assign each cell a
  majority fuel class
\item
  Statistics of the fuel raster within ecoregion polygons (raster
  clipped by polygon)
\end{itemize}

\hypertarget{attribute-filtering-and-aggregation}{%
\subsubsection{Attribute filtering and
aggregation}\label{attribute-filtering-and-aggregation}}

\begin{itemize}
\tightlist
\item
  Filtering fire records to NWCG\_CAUSE == ``Natural'' and to the date
  range February 2019--December 2020
\item
  Lightning will use the combined Monthly Thunder Hour product for 2019
  and 2020 to summ across CONUS to get a per grid cell count

  \begin{itemize}
  \tightlist
  \item
    Monthly Thunder Hour: did lightning happen in this cell during this
    hour, yes or no?
  \end{itemize}
\item
  Aggregating fuel pixel counts within each ecoregion into classes
\end{itemize}

\hypertarget{choropleth-mapping-and-advanced-visualization}{%
\subsubsection{Choropleth mapping and advanced
visualization}\label{choropleth-mapping-and-advanced-visualization}}

\begin{itemize}
\tightlist
\item
  Fuel map colored by their ignition-to-thunder-hour ratio
\item
  Ecoregions colored by their ignition-to-thunder-hour ratio
\item
  Create a map that displays the areas with the most lightning cause
  wildfire ignitions
\end{itemize}

    \includegraphics{/home/jovyan/work/Project2/Project2_OverviewofData.png}

    \begin{tcolorbox}[breakable, size=fbox, boxrule=1pt, pad at break*=1mm,colback=cellbackground, colframe=cellborder]
\prompt{In}{incolor}{1}{\boxspacing}
\begin{Verbatim}[commandchars=\\\{\}]
\PY{c+c1}{\PYZsh{} \PYZpc{}pip install earthaccess}
\PY{c+c1}{\PYZsh{} import earthaccess}
\PY{c+c1}{\PYZsh{}import glob}

\PY{c+c1}{\PYZsh{}files = sorted(glob.glob(\PYZdq{}data/glmmth/GLMc\PYZus{}th\PYZus{}*.nc\PYZdq{}))}
\PY{c+c1}{\PYZsh{} earthaccess.login(strategy=\PYZdq{}interactive\PYZdq{}, persist=True)}

\PY{c+c1}{\PYZsh{} results = earthaccess.search\PYZus{}data(}
\PY{c+c1}{\PYZsh{}     short\PYZus{}name=\PYZdq{}glmmth\PYZdq{},   }
\PY{c+c1}{\PYZsh{}     temporal=(\PYZdq{}2019\PYZhy{}01\PYZhy{}01\PYZdq{}, \PYZdq{}2020\PYZhy{}12\PYZhy{}31\PYZdq{}),}
\PY{c+c1}{\PYZsh{} )}
\PY{c+c1}{\PYZsh{} files = earthaccess.download(results, \PYZdq{}work/Project2/data\PYZdq{})}
\end{Verbatim}
\end{tcolorbox}

    \begin{tcolorbox}[breakable, size=fbox, boxrule=1pt, pad at break*=1mm,colback=cellbackground, colframe=cellborder]
\prompt{In}{incolor}{2}{\boxspacing}
\begin{Verbatim}[commandchars=\\\{\}]
\PY{c+c1}{\PYZsh{} For each file, promotes lat/lon to real coordinates and crops to CONUS, then sums each year\PYZsq{}s 12 monthly slices and adds both years together into }
\PY{c+c1}{\PYZsh{} one cumulative thunder\PYZhy{}hours grid}

\PY{c+c1}{\PYZsh{} Gridded because \PYZdq{}thunder hour\PYZdq{} is a binning concept by definition, and because it merges detections from instruments with different viewing geometries }
\PY{c+c1}{\PYZsh{} (both satellites, and more broadly ground networks) onto one common reference frame}

\PY{c+c1}{\PYZsh{} lat is 520 values, lon is 1180 values, and total is 1180×520 — like a spreadsheet: lat gives the row headers, lon gives the column headers, and cell }
\PY{c+c1}{\PYZsh{} total[i, j] is the thunder\PYZhy{}hour count}

\PY{k+kn}{import} \PY{n+nn}{glob}
\PY{k+kn}{import} \PY{n+nn}{xarray} \PY{k}{as} \PY{n+nn}{xr}

\PY{n}{files} \PY{o}{=} \PY{n+nb}{sorted}\PY{p}{(}\PY{n}{glob}\PY{o}{.}\PY{n}{glob}\PY{p}{(}\PY{l+s+s2}{\PYZdq{}}\PY{l+s+s2}{data/glmmth/GLMc\PYZus{}th\PYZus{}*.nc}\PY{l+s+s2}{\PYZdq{}}\PY{p}{)}\PY{p}{)}

\PY{k}{def} \PY{n+nf}{load\PYZus{}and\PYZus{}crop}\PY{p}{(}\PY{n}{path}\PY{p}{)}\PY{p}{:}
    \PY{n}{ds} \PY{o}{=} \PY{n}{xr}\PY{o}{.}\PY{n}{open\PYZus{}dataset}\PY{p}{(}\PY{n}{path}\PY{p}{)}
    \PY{c+c1}{\PYZsh{} lat/lon are plain data variables here, not indexed coordinates —}
    \PY{c+c1}{\PYZsh{} promote them so .sel() can crop by actual lat/lon values}
    \PY{n}{ds} \PY{o}{=} \PY{n}{ds}\PY{o}{.}\PY{n}{assign\PYZus{}coords}\PY{p}{(}\PY{n}{lat}\PY{o}{=}\PY{p}{(}\PY{l+s+s2}{\PYZdq{}}\PY{l+s+s2}{nlat}\PY{l+s+s2}{\PYZdq{}}\PY{p}{,} \PY{n}{ds}\PY{o}{.}\PY{n}{lat}\PY{o}{.}\PY{n}{values}\PY{p}{)}\PY{p}{,} \PY{n}{lon}\PY{o}{=}\PY{p}{(}\PY{l+s+s2}{\PYZdq{}}\PY{l+s+s2}{nlon}\PY{l+s+s2}{\PYZdq{}}\PY{p}{,} \PY{n}{ds}\PY{o}{.}\PY{n}{lon}\PY{o}{.}\PY{n}{values}\PY{p}{)}\PY{p}{)}
    \PY{n}{ds} \PY{o}{=} \PY{n}{ds}\PY{o}{.}\PY{n}{swap\PYZus{}dims}\PY{p}{(}\PY{p}{\PYZob{}}\PY{l+s+s2}{\PYZdq{}}\PY{l+s+s2}{nlat}\PY{l+s+s2}{\PYZdq{}}\PY{p}{:} \PY{l+s+s2}{\PYZdq{}}\PY{l+s+s2}{lat}\PY{l+s+s2}{\PYZdq{}}\PY{p}{,} \PY{l+s+s2}{\PYZdq{}}\PY{l+s+s2}{nlon}\PY{l+s+s2}{\PYZdq{}}\PY{p}{:} \PY{l+s+s2}{\PYZdq{}}\PY{l+s+s2}{lon}\PY{l+s+s2}{\PYZdq{}}\PY{p}{\PYZcb{}}\PY{p}{)}

    \PY{c+c1}{\PYZsh{} handle either latitude ordering (ascending or descending)}
    \PY{n}{lat\PYZus{}vals} \PY{o}{=} \PY{n}{ds}\PY{o}{.}\PY{n}{lat}\PY{o}{.}\PY{n}{values}
    \PY{n}{lat\PYZus{}slice} \PY{o}{=} \PY{n+nb}{slice}\PY{p}{(}\PY{l+m+mi}{24}\PY{p}{,} \PY{l+m+mi}{50}\PY{p}{)} \PY{k}{if} \PY{n}{lat\PYZus{}vals}\PY{p}{[}\PY{l+m+mi}{0}\PY{p}{]} \PY{o}{\PYZlt{}} \PY{n}{lat\PYZus{}vals}\PY{p}{[}\PY{o}{\PYZhy{}}\PY{l+m+mi}{1}\PY{p}{]} \PY{k}{else} \PY{n+nb}{slice}\PY{p}{(}\PY{l+m+mi}{50}\PY{p}{,} \PY{l+m+mi}{24}\PY{p}{)}

    \PY{k}{return} \PY{n}{ds}\PY{o}{.}\PY{n}{sel}\PY{p}{(}\PY{n}{lat}\PY{o}{=}\PY{n}{lat\PYZus{}slice}\PY{p}{,} \PY{n}{lon}\PY{o}{=}\PY{n+nb}{slice}\PY{p}{(}\PY{o}{\PYZhy{}}\PY{l+m+mi}{125}\PY{p}{,} \PY{o}{\PYZhy{}}\PY{l+m+mi}{66}\PY{p}{)}\PY{p}{)}

\PY{n}{datasets} \PY{o}{=} \PY{p}{[}\PY{n}{load\PYZus{}and\PYZus{}crop}\PY{p}{(}\PY{n}{f}\PY{p}{)} \PY{k}{for} \PY{n}{f} \PY{o+ow}{in} \PY{n}{files}\PY{p}{]}

\PY{c+c1}{\PYZsh{} If range assumptions above don\PYZsq{}t match this file and need adjusting}
\PY{n+nb}{print}\PY{p}{(}\PY{l+s+s2}{\PYZdq{}}\PY{l+s+s2}{Cropped shape:}\PY{l+s+s2}{\PYZdq{}}\PY{p}{,} \PY{n}{datasets}\PY{p}{[}\PY{l+m+mi}{0}\PY{p}{]}\PY{o}{.}\PY{n}{thunder\PYZus{}hours}\PY{o}{.}\PY{n}{shape}\PY{p}{)}

\PY{c+c1}{\PYZsh{} sum each file\PYZsq{}s 12 months, then sum both years together}
\PY{n}{total} \PY{o}{=} \PY{n+nb}{sum}\PY{p}{(}\PY{n}{ds}\PY{o}{.}\PY{n}{thunder\PYZus{}hours}\PY{o}{.}\PY{n}{sum}\PY{p}{(}\PY{n}{dim}\PY{o}{=}\PY{l+s+s2}{\PYZdq{}}\PY{l+s+s2}{nmon}\PY{l+s+s2}{\PYZdq{}}\PY{p}{)} \PY{k}{for} \PY{n}{ds} \PY{o+ow}{in} \PY{n}{datasets}\PY{p}{)}

\PY{n+nb}{print}\PY{p}{(}\PY{n}{total}\PY{p}{)}
\end{Verbatim}
\end{tcolorbox}

    \begin{Verbatim}[commandchars=\\\{\}]
Cropped shape: (1180, 520, 12)
<xarray.DataArray 'thunder\_hours' (lon: 1180, lat: 520)>
array([[ 11,  10,   9, {\ldots},   0,   0,   0],
       [ 16,  14,  12, {\ldots},   0,   0,   0],
       [ 21,  17,  16, {\ldots},   0,   0,   0],
       {\ldots},
       [ 14,  16,  17, {\ldots}, 184, 181, 182],
       [ 19,  17,  16, {\ldots}, 180, 185, 197],
       [ 20,  18,  16, {\ldots}, 183, 191, 195]])
Coordinates:
  * lat      (lat) float32 49.97 49.92 49.88 49.83 {\ldots} 24.17 24.12 24.08 24.02
  * lon      (lon) float32 -125.0 -124.9 -124.9 -124.8 {\ldots} -66.12 -66.07 -66.03
    \end{Verbatim}

    \begin{tcolorbox}[breakable, size=fbox, boxrule=1pt, pad at break*=1mm,colback=cellbackground, colframe=cellborder]
\prompt{In}{incolor}{3}{\boxspacing}
\begin{Verbatim}[commandchars=\\\{\}]
\PY{c+c1}{\PYZsh{} counts of each fuel class per polygon}
\PY{k+kn}{import} \PY{n+nn}{os}
\PY{k+kn}{import} \PY{n+nn}{geopandas} \PY{k}{as} \PY{n+nn}{gpd}
\PY{k+kn}{import} \PY{n+nn}{pandas} \PY{k}{as} \PY{n+nn}{pd}

\PY{n}{cache\PYZus{}path} \PY{o}{=} \PY{l+s+s2}{\PYZdq{}}\PY{l+s+s2}{data/processed/fire\PYZus{}natural\PYZus{}2019\PYZus{}2020.gpkg}\PY{l+s+s2}{\PYZdq{}}

\PY{k}{if} \PY{n}{os}\PY{o}{.}\PY{n}{path}\PY{o}{.}\PY{n}{exists}\PY{p}{(}\PY{n}{cache\PYZus{}path}\PY{p}{)}\PY{p}{:}
    \PY{n}{fire} \PY{o}{=} \PY{n}{gpd}\PY{o}{.}\PY{n}{read\PYZus{}file}\PY{p}{(}\PY{n}{cache\PYZus{}path}\PY{p}{)}
\PY{k}{else}\PY{p}{:}
    \PY{n}{fire} \PY{o}{=} \PY{n}{gpd}\PY{o}{.}\PY{n}{read\PYZus{}file}\PY{p}{(}\PY{l+s+s2}{\PYZdq{}}\PY{l+s+s2}{data/processed/fire\PYZus{}occurrence\PYZus{}albers.gpkg}\PY{l+s+s2}{\PYZdq{}}\PY{p}{)}

    \PY{n}{fire}\PY{p}{[}\PY{l+s+s2}{\PYZdq{}}\PY{l+s+s2}{DISCOVERY\PYZus{}DATE}\PY{l+s+s2}{\PYZdq{}}\PY{p}{]} \PY{o}{=} \PY{n}{pd}\PY{o}{.}\PY{n}{to\PYZus{}datetime}\PY{p}{(}\PY{n}{fire}\PY{p}{[}\PY{l+s+s2}{\PYZdq{}}\PY{l+s+s2}{DISCOVERY\PYZus{}DATE}\PY{l+s+s2}{\PYZdq{}}\PY{p}{]}\PY{p}{)}
    \PY{n}{fire} \PY{o}{=} \PY{n}{fire}\PY{p}{[}
        \PY{p}{(}\PY{n}{fire}\PY{p}{[}\PY{l+s+s2}{\PYZdq{}}\PY{l+s+s2}{NWCG\PYZus{}CAUSE\PYZus{}CLASSIFICATION}\PY{l+s+s2}{\PYZdq{}}\PY{p}{]} \PY{o}{==} \PY{l+s+s2}{\PYZdq{}}\PY{l+s+s2}{Natural}\PY{l+s+s2}{\PYZdq{}}\PY{p}{)}
        \PY{o}{\PYZam{}} \PY{p}{(}\PY{n}{fire}\PY{p}{[}\PY{l+s+s2}{\PYZdq{}}\PY{l+s+s2}{DISCOVERY\PYZus{}DATE}\PY{l+s+s2}{\PYZdq{}}\PY{p}{]} \PY{o}{\PYZgt{}}\PY{o}{=} \PY{l+s+s2}{\PYZdq{}}\PY{l+s+s2}{2019\PYZhy{}02\PYZhy{}12}\PY{l+s+s2}{\PYZdq{}}\PY{p}{)}
        \PY{o}{\PYZam{}} \PY{p}{(}\PY{n}{fire}\PY{p}{[}\PY{l+s+s2}{\PYZdq{}}\PY{l+s+s2}{DISCOVERY\PYZus{}DATE}\PY{l+s+s2}{\PYZdq{}}\PY{p}{]} \PY{o}{\PYZlt{}}\PY{o}{=} \PY{l+s+s2}{\PYZdq{}}\PY{l+s+s2}{2020\PYZhy{}12\PYZhy{}31}\PY{l+s+s2}{\PYZdq{}}\PY{p}{)}
    \PY{p}{]}

    \PY{n}{fire}\PY{o}{.}\PY{n}{to\PYZus{}file}\PY{p}{(}\PY{n}{cache\PYZus{}path}\PY{p}{,} \PY{n}{driver}\PY{o}{=}\PY{l+s+s2}{\PYZdq{}}\PY{l+s+s2}{GPKG}\PY{l+s+s2}{\PYZdq{}}\PY{p}{)}
\end{Verbatim}
\end{tcolorbox}

    \begin{tcolorbox}[breakable, size=fbox, boxrule=1pt, pad at break*=1mm,colback=cellbackground, colframe=cellborder]
\prompt{In}{incolor}{4}{\boxspacing}
\begin{Verbatim}[commandchars=\\\{\}]
\PY{n+nb}{print}\PY{p}{(}\PY{n}{fire}\PY{o}{.}\PY{n}{head}\PY{p}{(}\PY{p}{)}\PY{p}{)}
\end{Verbatim}
\end{tcolorbox}

    \begin{Verbatim}[commandchars=\\\{\}]
         OBJECTID     FOD\_ID      FPA\_ID SOURCE\_SYSTEM\_TYPE SOURCE\_SYSTEM  \textbackslash{}
2166748   2166749  400500002  FS-6955822                FED   FS-FIRESTAT
2166749   2166750  400500003  FS-6936597                FED   FS-FIRESTAT
2166751   2166752  400500005  FS-6955006                FED   FS-FIRESTAT
2166752   2166753  400500006  FS-6947814                FED   FS-FIRESTAT
2166753   2166754  400500007  FS-6957509                FED   FS-FIRESTAT

        NWCG\_REPORTING\_AGENCY NWCG\_REPORTING\_UNIT\_ID  \textbackslash{}
2166748                    FS                USMTKNF
2166749                    FS                USMTKNF
2166751                    FS                USMTKNF
2166752                    FS                USMTKNF
2166753                    FS                USMTHLF

                         NWCG\_REPORTING\_UNIT\_NAME SOURCE\_REPORTING\_UNIT  \textbackslash{}
2166748                  Kootenai National Forest                  0114
2166749                  Kootenai National Forest                  0114
2166751                  Kootenai National Forest                  0114
2166752                  Kootenai National Forest                  0114
2166753  Helena - Lewis and Clark National Forest                  0115

                       SOURCE\_REPORTING\_UNIT\_NAME  {\ldots} FIRE\_SIZE  \textbackslash{}
2166748                  Kootenai National Forest  {\ldots}       0.1
2166749                  Kootenai National Forest  {\ldots}       0.6
2166751                  Kootenai National Forest  {\ldots}       0.1
2166752                  Kootenai National Forest  {\ldots}       0.1
2166753  Helena - Lewis and Clark National Forest  {\ldots}       0.1

        FIRE\_SIZE\_CLASS   LATITUDE   LONGITUDE OWNER\_DESCR STATE COUNTY  \textbackslash{}
2166748               A  48.882222 -115.020556     PRIVATE    MT    053
2166749               B  48.484167 -115.258056        USFS    MT    053
2166751               A  48.457500 -115.913333        USFS    MT    053
2166752               A  48.647778 -114.830556        USFS    MT    029
2166753               A  46.913333 -112.488889        USFS    MT    049

        FIPS\_CODE               FIPS\_NAME                          geometry
2166748     30053          Lincoln County  POINT (-1402356.177 3016090.463)
2166749     30053          Lincoln County  POINT (-1428438.760 2976734.126)
2166751     30053          Lincoln County  POINT (-1476952.178 2983877.638)
2166752     30029         Flathead County  POINT (-1393606.561 2988029.143)
2166753     30049  Lewis and Clark County  POINT (-1255217.511 2767266.359)

[5 rows x 39 columns]
    \end{Verbatim}

    \begin{Verbatim}[commandchars=\\\{\}, frame=single, framerule=2mm, rulecolor=\color{outerrorbackground}]
\textcolor{ansi-red}{---------------------------------------------------------------------------}
\textcolor{ansi-red}{NameError}                                 Traceback (most recent call last)
Cell \textcolor{ansi-green}{In[4], line 2}
\textcolor{ansi-green-intense}{\textbf{      1}} \def\tcRGB{\textcolor[RGB]}\expandafter\tcRGB\expandafter{\detokenize{0,135,0}}{print}(fire\def\tcRGB{\textcolor[RGB]}\expandafter\tcRGB\expandafter{\detokenize{98,98,98}}{.}head())
\textcolor{ansi-green}{----> 2} \def\tcRGB{\textcolor[RGB]}\expandafter\tcRGB\expandafter{\detokenize{0,135,0}}{print}(\setlength{\fboxsep}{0pt}\colorbox{ansi-yellow}{ecoregions\strut}\def\tcRGB{\textcolor[RGB]}\expandafter\tcRGB\expandafter{\detokenize{98,98,98}}{.}head())

\textcolor{ansi-red}{NameError}: name 'ecoregions' is not defined
    \end{Verbatim}


    % Add a bibliography block to the postdoc
    
    
    
\end{document}
